\documentclass[10pt,a4paper,notitlepage,twocolumn]{article}
\usepackage{amsmath,amssymb,amsbsy}
\usepackage{float}
\usepackage[french]{babel}
\usepackage{graphicx}

\usepackage[utf8x]{inputenc}
\usepackage[T1]{fontenc}
\usepackage{palatino}
\usepackage{manfnt}
\usepackage{hyperref}
\usepackage{nicefrac}

\usepackage[active]{srcltx}
\usepackage{scrtime}

\newcommand{\exercice}[1]{\textsc{\textbf{Exercice}} #1}
\newcommand{\question}[1]{\textbf{(#1)}}
\setlength{\parindent}{0cm}

\begin{document}

\title{\textsc{Croissance\\ \small{(L2, Le Mans)}}}
\author{Stéphane Adjemian\thanks{Université du Mans. \texttt{stephane
DOT adjemian AT univ DASH lemans DOT fr}}}
\date{Mercredi 18 décembre 2024}

\maketitle

\thispagestyle{empty}

\begin{quote}
  \textit{Les réponses non commentées ou insuffisamment détaillées ne
  seront pas considérées. Prenez le temps d'écrire des phrases.}
\end{quote}

\bigskip
\bigskip

On considère un modèle de Solow augmenté de dépenses publiques
productives (dans le sens où la dépense de l'état et nécessaire et
accroît la production). Soit la fonction de production~:
\[
Y(t) = K(t)^{\alpha}G(t)^{\beta}(L(t))^{1-\alpha-\beta}
\]
Les variables $Y$, $K$ et $L$ ont les interprétations usuelles. La
population, $L$, croît au taux constant $n>0$. La variable $G$ est le
niveau de la dépense publique.  Les paramètres $\alpha$ et $\beta$
sont positifs et vérifient $\alpha+\beta<1$. La dépense publique est
financée par une taxe, $0<\tau<1$, sur la
production $G(t) = \tau Y(t)$. On note $0<s<1$ le taux
d'épargne. L'investissement à l'instant $t$ est donc donné
par $I(t) = s (1-\tau) Y(t)$ et la consommation à l'instant $t$
par $C(t) = (1-s)(1-\tau)Y(t)$.\newline

\question{1} Montrer que la dynamique du stock de capital par tête est caractérisée par~:

\[
\dot{k} = s(1-\tau)k^{\alpha}g^{\beta} - (\delta+n)k
\]
où $\delta>0$ est le taux de dépréciation du stock de capital physique.\newline

\question{2} En notant que l'on peut exprimer la dépense publique par
tête comme une fonction linéaire de la production par tête $g=\tau y$,
montrer que l'on a aussi~:
\[
\dot{k} = s(1-\tau)\tau^{\frac{\beta}{1-\beta}}k^{\frac{\alpha}{1-\beta}} - (n+\delta)k
\]
Interpréter la forme de cette équation.\newline

\question{3} Donner une expression analytique du taux de croissance du
stock de capital par tête. Représenter graphiquement ce taux de
croissance.\newline

\question{4} Montrer l'existence et calculer l'état stationnaire non
trivial pour le stock de capital par tête, on notera $k^{\star}$. En
déduire l'état stationnaire pour la production par tête et la
consommation par tête, on notera $y^\star$ et $c^{\star}$.\newline

\question{5} Comment se comporte les variables par tête à long terme~?\newline

\question{6} Expliquer pourquoi nous excluons $\tau=0$ et $\tau=1$.\newline

\question{7} Calculer la taxe optimale, notée $\tau^{\star}$ qui
maximise la consommation par tête à long terme.\newline

\question{8} Dans le choix de la taxation optimale, existe t-il un
arbitrage entre niveau de la consommation à long terme et croissance
des variables par tête~?\newline

\question{9} En procédant par analogie avec les résultats obenus en
cours, c'est-à-dire sans refaire les calculs, déterminer la vitesse
d'ajustement vers l'état stationnaire (nous noterons $\lambda$ la
vitesse de convergence). Commenter.\newline

\question{10} Que deviennent les prédictions du modèle quand $\beta=1-\alpha$~?

\end{document}
